\subsection{QC-005: an explicit surviving degree-six highest-weight copy}

\paragraph{Statement QC-005.}
There is an explicit degree-six highest-weight polynomial of weight
$\lambda=(15,13,13,13)$ on quaternary nonics that is nonzero on
the variety $V=\{\ell q:\deg q=8\}$. Thus $q_\lambda\geq1$.
Combined with the ambient-multiplicity-one calculation in the
main packet, this gives $q_\lambda=1$.
This statement does not assert that it is an equation of $D_{9,4}$.

\begin{proof}
Introduce six four-component variable vectors $y_0,\ldots,y_5$.
For $i<j<k<l$, let
\[
 \Delta_{ijkl}=
 \det\begin{pmatrix}
 \partial_{y_{i0}}&\partial_{y_{j0}}&\partial_{y_{k0}}&\partial_{y_{l0}}\\
 \partial_{y_{i1}}&\partial_{y_{j1}}&\partial_{y_{k1}}&\partial_{y_{l1}}\\
 \partial_{y_{i2}}&\partial_{y_{j2}}&\partial_{y_{k2}}&\partial_{y_{l2}}\\
 \partial_{y_{i3}}&\partial_{y_{j3}}&\partial_{y_{k3}}&\partial_{y_{l3}}
 \end{pmatrix}.
\]
Define the constant-coefficient differential operator
\[
 \Omega=\Delta_{0135}^{4}\Delta_{1235}
                \Delta_{2345}^{3}\Delta_{0234}\Delta_{0124}^{4}.
 \tag{QC-005.1}
\]
For a nonic $F$, apply it to $\prod_{i=0}^5F(y_i)$.
The orders in the six copies are $(9,9,9,9,8,8)$, so the result
is constant in the first four vectors and bilinear in the last
two. Set $y_4=y_5=e_0=(1,0,0,0)$ and call the result $P(F)$.
It is an integral polynomial of degree six in the ordinary
coefficients of $F$.

To check its weight and highest-weight property without a
differential convention ambiguity, polarize using formal vectors
$v_0,\ldots,v_5$. Apart from the constant $(9!)^6$, its symbol is
\[
 [0135]^4[1235][2345]^3[0234][0124]^4\,(v_4)_0(v_5)_0,
 \tag{QC-005.2}
\]
where a bracket is the determinant of the four indicated vectors.
Each letter has degree nine. Passing to the symmetric polarization
in the six letters produces an element of
$\operatorname{Sym}^6(\operatorname{Sym}^9\mathbb C^4)$ in the
convention of the prompt. Every bracket contributes weight
$(1,1,1,1)$, and the two free factors contribute $(2,0,0,0)$.
There are thirteen brackets, giving weight $(15,13,13,13)$.
All positive-root operators kill a determinant bracket and kill
the free highest-weight coordinate, so they kill (QC-005.2) and
its symmetric polarization. It remains to prove it is nonzero.

Use the following integer data, with rows $(b_i,a_{i1},a_{i2},a_{i3},w_i)$:
\[
\begin{pmatrix}
90&69&89&83&81\\
98&94&54&38&33\\
95&46&16&17&77\\
24&52&62&65&44
\end{pmatrix}.
\]
Put $a_i=a_{i1}x_1+a_{i2}x_2+a_{i3}x_3$ and
\[
 F_* = \sum_{i=1}^4w_i\bigl((b_ix_0+a_i)^9-a_i^9\bigr).
 \tag{QC-005.3}
\]
Every summand is divisible by $x_0$ over $\mathbb Z$, so
$F_*=x_0q_*$ for an octic $q_*$ and $F_*\in V$ exactly.

The standalone checker \texttt{check\_nonic\_candidate.py}
evaluates (QC-005.2) on this eight-term sum of ninth powers.
Every pair of symbolic letters lies together in at least one
bracket. Thus an assignment repeating a vector has zero bracket,
and only the $8\cdot7\cdot6\cdot5\cdot4\cdot3=20160$
injective assignments need to be summed. For each assignment,
multiply the thirteen determinant values, the six power-sum
coefficients, and the two indicated free coordinates. This is
exactly multilinearity of (QC-005.1), with no stochastic test or
assumption about generic points. Direct integer arithmetic modulo
$101$ gives
\[
 (9!)^{-6}P(F_*)=39\pmod{101},\qquad
 P(F_*)=34\pmod{101}.
 \tag{QC-005.4}
\]
The factor $(9!)^6$ comes from differentiating each ninth power
eight or nine times; the remaining degree is one or zero, so
both cases have the factor $9!$. Since $101>9$, division in the
first expression is legitimate. The second expression proves
that the integral polynomial is nonzero at an integer point of
$V$, hence nonzero over $\mathbb C$.

By the settled CUE\_GCT R-13 at $m=8$, the four-variable octic
permanent image is the whole octic space. The four-variable
reduction therefore transfers this nonzero copy to the padded
permanent orbit closure. This proves $q_\lambda\geq1$, and the
separately checked ambient upper bound one gives equality.
\end{proof}

\paragraph{The determinant test, resolved in QC-006.}
The nonzero padded evaluation is only one side of the comparison.
QC-006 supplies a nonzero determinant evaluation of this same polynomial,
so the complete multiplicity comparison is $K_9(\lambda)=q_\lambda=1$.
The symmetric rectangular Kronecker upper bound is much larger than one
and does not decide the comparison by itself.

\paragraph{Geometric restatement.}
Equation (QC-005.1) is an integral Schur-functor construction;
after extension to characteristic zero its highest-weight line
generates the indicated irreducible equivariant submodule.
Equation (QC-005.4) says that the restriction of this section to
the coherent coordinate sheaf of the linear-factor locus is
nonzero. QC-006 separately proves that this section does not belong to the
equivariant ideal sheaf of the determinantal locus.
The framework locators are the SGA 3 and SGA 6 weight and
Schur-operation dictionary specified above; no purity or motive
statement is needed.
